% Source: https://artofproblemsolving.com/community/c3h2297782_easy_mathz
% Stable identity: aops-topic-2297782-q01
Multiply the equation
\[\dfrac{x+y+z}{3}=\dfrac{\dfrac{x+y}{2}+z}{2}\]
by $12$: $4(x+y+z)=3(x+y)+6z$, that is, $x+y=2z$. So the two methods agree exactly when
\[\boxed{z=\dfrac{x+y}{2}},\]
that is, when the third score equals the average of the first two (and then both methods give $z$). Otherwise they differ, because the second method gives weight $\dfrac12$ to $z$ and only $\dfrac14$ to each of $x$ and $y$, instead of $\dfrac13$ to each.
